%=========================================
% Section: Model Architecture
% Drop-in fragment — paste into your report
%
% Preamble requirements (add once):
%   \usepackage{tikz}
%   \usetikzlibrary{arrows.meta,backgrounds,calc}
%   \usepackage{amsmath,amssymb}
%   \usepackage{booktabs,array}
%=========================================

\section{Model Architecture}
\label{sec:model_arch}

\noindent
This section presents the complete architecture of the proposed four-objective
Electric Vehicle Charging Station (EVCS) optimization framework designed for
Indore city. The framework integrates five tightly coupled components in a
unified, data-driven pipeline:

\begin{itemize}
    \item a logistic demand forecasting and Monte~Carlo stochastic simulation
          module,
    \item a four-objective binary optimisation layer using NSGA-II and MOPSO,
    \item an M/M/$c$/K finite-capacity queueing evaluation layer with
          Time-of-Use demand segmentation,
    \item a blocking probability constraint enforcement mechanism, and
    \item a sensitivity analysis and scalability benchmarking module.
\end{itemize}

\noindent
These components interact through shared spatial datasets spanning the 85
administrative wards of Indore, a heterogeneous pool of $30+$ candidate venues
across five functional zones, and a logistic EV adoption curve anchored to
CAGR-based ward-level demographic projections. The result is a scalable
framework capable of generating Pareto-optimal EVCS deployment strategies
evaluated against formal quality indicators: Hypervolume~(HV), Generational
Distance~(GD), and Spread.

\begin{center}
\fbox{\parbox{0.93\linewidth}{
\textbf{Data Layer}:\ \ Census 2011 $\rightarrow$ CAGR $\rightarrow$ Logistic
Adoption Curve $\rightarrow$ Ward-level Daily Demand \\[4pt]
\textbf{Stochastic Layer}:\ \ Monte~Carlo Simulation (100 scenarios,
$v = 0.15$) \\[4pt]
\textbf{Optimization Layer}:\ \ NSGA-II $\|$ MOPSO
$\rightarrow$ Pareto-Optimal Station Subsets \\[4pt]
\textbf{Evaluation Layer}:\ \ M/M/$c$/K Queue (ToU) + $P_{\text{block}}$
Penalty + HV / GD / Spread \\[4pt]
\textbf{Output}:\ \ Convergence Plots + Pareto Fronts + Sensitivity Analysis
+ Deployment Recommendation
}}
\end{center}



%-----------------------------------------
\section{Baseline Model}
\label{sec:baseline_model}

\noindent
The baseline model provides a deterministic, single-period, cost--coverage
formulation used as a reference benchmark for evaluating the multi-objective
optimization results. It focuses exclusively on minimizing infrastructure cost
while ensuring citywide spatial coverage under grid-capacity constraints.

\subsection{Deterministic Cost--Coverage Formulation}

Let
\begin{align*}
\mathcal{W} &:\ \text{set of demand wards},         & |\mathcal{W}| &= 85,\\
\mathcal{J} &:\ \text{set of candidate sites},      & |\mathcal{J}| &= 30+,\\
\mathcal{Z} &:\ \text{set of functional zones},     & |\mathcal{Z}| &= 5.
\end{align*}

\noindent
\textbf{Decision variable:}
\begin{equation}
x_j =
\begin{cases}
1, & \text{if candidate station } j \text{ is selected for deployment,}\\
0, & \text{otherwise,}
\end{cases}
\qquad j \in \mathcal{J}.
\label{eq:decision}
\end{equation}
\textit{where:} $x_j \in \{0,1\}$ is the binary selection variable for
candidate site $j$; $\mathcal{J}$ is the full set of candidate EVCS sites.

\medskip\noindent
The deterministic objective minimizes total capital expenditure:
\begin{equation}
\mathrm{CAPEX}
= \sum_{j \in \mathcal{J}} x_j
  \bigl(C_{\mathrm{base},j} + C_{\mathrm{grid},j}\bigr),
\label{eq:capex}
\end{equation}
\textit{where:}
\begin{itemize}
  \item $C_{\mathrm{base},j}$ — venue-type-dependent base opening cost
        (\rupee6.5--18.0~lakh per Table~3.2),
  \item $C_{\mathrm{grid},j}$ — electrical grid upgrade cost at site $j$.
\end{itemize}

\medskip\noindent
Coverage constraints ensure every ward is within the zone-adjusted
service radius $R_z$:
\begin{equation}
\sum_{j \in \mathcal{J}} x_j \,\mathbb{1}
\!\left(\mathrm{dist}_{wj} \leq R_z\right) \geq 1,
\quad \forall\, w \in \mathcal{W},
\label{eq:coverage}
\end{equation}
\textit{where:}
\begin{itemize}
  \item $\mathrm{dist}_{wj}$ — Euclidean distance from ward $w$ centroid
        to site $j$ (km),
  \item $\mathbb{1}(\cdot)$ — indicator function; equals 1 if
        $\mathrm{dist}_{wj}\leq R_z$, else 0,
  \item $R_z = 1.5$~km (Central Core) and $R_z = 2.5$~km (Outer zones),
        per MoRTH~(2022).
\end{itemize}

\subsection{Baseline Architecture Interpretation}

The deterministic model establishes a cost-minimizing reference network that
covers all 85~wards. Although useful as a lower-bound benchmark on capital
expenditure, this baseline does not capture:

\begin{itemize}
    \item stochastic EV demand variability and Monte~Carlo uncertainty,
    \item multi-objective trade-offs among cost, coverage, ROI, and
          service quality,
    \item financial viability (net return on investment),
    \item queue waiting times or blocking probability under peak demand,
    \item the distinction between NSGA-II and MOPSO solution quality.
\end{itemize}

\noindent
These limitations motivate the extended four-objective hybrid framework
described in the following section.

%-----------------------------------------
\section{Proposed Four-Objective Optimization Framework}
\label{sec:hybrid_framework}

\noindent
The complete framework jointly optimises station placement across four
conflicting objectives, enforces queue service constraints, and evaluates
robustness under stochastic demand. Its five interacting components are
described below.

%-----------------------------------------
\subsection{Demand Forecasting and Monte~Carlo Simulation Layer}
\label{sec:demand_layer}

Ward-level EV charging demand is projected through a two-stage pipeline.
First, the population of ward $w$ at forecast year $t$ is computed via a
constant CAGR:
\begin{equation}
\mathrm{Pop}_w(t)
= \mathrm{Pop}_w(2011)\times(1 + r)^{\,t-2011},
\qquad r = 0.018.
\label{eq:cagr}
\end{equation}
\textit{where:}
\begin{itemize}
  \item $\mathrm{Pop}_w(2011)$ — Census 2011 base population of ward $w$,
  \item $r = 0.018$ — annual compound growth rate (CAGR),
  \item $t$ — forecast year.
\end{itemize}

\medskip\noindent
Second, the aggregate EV penetration rate follows a logistic S-curve:
\begin{equation}
E(t) = \frac{K}{1 + \exp\!\bigl(-a\,(t - t_0)\bigr)},
\qquad K=1.0,\; a=0.3,\; t_0=2030.
\label{eq:scurve}
\end{equation}
\textit{where:}
\begin{itemize}
  \item $K = 1.0$ — carrying capacity (maximum EV penetration level),
  \item $a = 0.3$ — steepness of the adoption growth curve,
  \item $t_0 = 2030$ — inflection year at which $E(t_0) = K/2$.
\end{itemize}

\medskip\noindent
This yields a daily public charging demand per ward:
\begin{equation}
D_w(t)
= \mathrm{Pop}_w(t)
  \times \frac{V}{1000}
  \times E(t)
  \times F_{\mathrm{public}},
\label{eq:demand}
\end{equation}
\textit{where:}
\begin{itemize}
  \item $V = 50$ — EV density (EVs per 1000 persons),
  \item $E(t)$ — EV penetration rate from Eq.~\eqref{eq:scurve},
  \item $F_{\mathrm{public}} = 0.20$ — fraction of EV owners relying on
        public charging infrastructure.
\end{itemize}

\medskip\noindent
To capture stochastic demand variability, 100~Monte~Carlo scenarios are
generated per candidate station $j$:
\begin{equation}
D_{j,s}
\sim \max\!\bigl(0,\;\mathcal{N}(\mu_j,\,(v\mu_j)^2)\bigr),
\qquad s = 1,\dots,100,\quad v = 0.15.
\label{eq:mc}
\end{equation}
\textit{where:}
\begin{itemize}
  \item $\mu_j$ — mean expected daily demand at station $j$,
  \item $v = 0.15$ — coefficient of variation for demand noise,
  \item $\mathcal{N}(\mu,\sigma^2)$ — normal distribution with mean $\mu$
        and variance $\sigma^2$,
  \item $\max(0,\cdot)$ — clamp enforcing non-negative demand values.
\end{itemize}
Financial and queueing objectives ($f_3$, $f_4$) are evaluated as averages
across all 100~scenarios, ensuring robustness against demand edge cases.

%-----------------------------------------
\subsection{Four-Objective Problem Formulation}
\label{sec:four_obj}

The EVCS placement problem is formulated as a mixed-integer nonlinear
programming (MINLP) model. Let $\mathbf{x} = [x_1,\dots,x_N]^{\top}$
be the binary decision vector. Objectives for maximisation are negated so
the solver uniformly minimises:
\begin{equation}
\mathbf{F}(\mathbf{x})
= \bigl[f_1(\mathbf{x}),\;f_2(\mathbf{x}),\;f_3(\mathbf{x}),\;
         f_4(\mathbf{x})\bigr]^{\!\top},
\label{eq:objvec}
\end{equation}
\textit{where:}
\begin{itemize}
  \item $\mathbf{x} \in \{0,1\}^N$ — binary station-selection vector,
  \item $N = |\mathcal{J}|$ — total number of candidate sites,
  \item $f_1,\ldots,f_4$ — the four scalar objectives defined below
        (Eqs.~\eqref{eq:f1}--\eqref{eq:f4}).
\end{itemize}

\subsubsection*{Objective $f_1$ — Total Capital Expenditure (Minimize)}
\begin{equation}
f_1(\mathbf{x})
= \sum_{j \in \mathcal{J}} x_j
  \bigl(C_{\mathrm{base},j} + C_{\mathrm{grid},j}\bigr).
\label{eq:f1}
\end{equation}
\textit{where:} $C_{\mathrm{base},j}$ — base infrastructure cost at site
$j$ (\rupee\ lakh); $C_{\mathrm{grid},j}$ — grid connection upgrade cost
at site $j$ (\rupee\ lakh).

\subsubsection*{Objective $f_2$ — Spatial Demand Coverage (Maximize)}
\begin{equation}
f_2(\mathbf{x})
= -\frac{\displaystyle\sum_{w \in \mathcal{W}}
    D_w \cdot \mathbb{1}\!\left(
    \min_{j:\,x_j=1}\mathrm{dist}_{wj} \leq R_z\right)}
  {\displaystyle\sum_{w \in \mathcal{W}} D_w}.
\label{eq:f2}
\end{equation}
\textit{where:}
\begin{itemize}
  \item $D_w$ — daily demand of ward $w$ from Eq.~\eqref{eq:demand},
  \item $\mathbb{1}(\cdot)$ — indicator: 1 if ward $w$ is within $R_z$
        of at least one selected station, else 0,
  \item $\mathrm{dist}_{wj}$ — distance from ward $w$ centroid to site
        $j$ (km); negation converts maximisation to minimisation.
\end{itemize}

\subsubsection*{Objective $f_3$ — Net Return on Investment (Maximize)}
\begin{equation}
f_3(\mathbf{x})
= -\frac{\text{Annual Revenue}(\mathbf{x})
         - \text{Annual OpEx}(\mathbf{x})}
        {f_1(\mathbf{x})}.
\label{eq:f3}
\end{equation}
\textit{where:}
\begin{itemize}
  \item Annual Revenue$(\mathbf{x})$ — total yearly charging revenue
        summed across selected stations (Monte~Carlo averaged),
  \item Annual OpEx$(\mathbf{x})$ — yearly operating expenditure
        (maintenance, electricity, staffing),
  \item $f_1(\mathbf{x})$ — CAPEX from Eq.~\eqref{eq:f1}; normalisation
        by CAPEX prevents near-perfect linear correlation with $f_2$
        ($r > 0.99$) that would degenerate the Pareto front.
\end{itemize}

\subsubsection*{Objective $f_4$ — Mean Queue Waiting Time (Minimize)}
\begin{equation}
f_4(\mathbf{x})
= \frac{1}{|\{j : x_j = 1\}|}
  \sum_{j:\,x_j=1} \overline{W}_{q,j}.
\label{eq:f4}
\end{equation}
\textit{where:}
\begin{itemize}
  \item $|\{j : x_j = 1\}|$ — number of deployed stations,
  \item $\overline{W}_{q,j}$ — ToU-weighted mean waiting time at station
        $j$, evaluated via the M/M/$c$/K model
        (Section~\ref{sec:queue_layer}),
  \item a penalty of $+50$~min is added to $f_4$ for every station
        where $P_{\mathrm{block}} > 10\%$.
\end{itemize}

%-----------------------------------------
\subsection{NSGA-II Optimization Layer}
\label{sec:nsgaii_layer}

NSGA-II is the primary evolutionary optimizer. Each individual in the
population encodes the binary station-selection vector
$\mathbf{x} \in \{0,1\}^N$. The algorithm proceeds as follows:

\begin{enumerate}
    \item \textbf{Initialisation:} A random population of size 80 is
          generated.
    \item \textbf{Evaluation:} Each individual is evaluated on all four
          objectives, including the M/M/$c$/K queueing model and
          Monte~Carlo demand averaging.
    \item \textbf{Non-dominated sorting:} Solutions are ranked into
          Pareto fronts $\mathcal{F}_1 \subset \mathcal{F}_2 \subset \cdots$
          using fast non-dominated sorting with
          $\mathcal{O}(MN^2)$ complexity.
    \item \textbf{Crowding distance selection:} Within the same front,
          solutions in less-crowded objective regions are preferred,
          maintaining spread diversity.
    \item \textbf{Variation:} Offspring are produced via uniform crossover
          ($p_c = 0.9$) and bit-flip mutation ($p_m = 0.08$).
    \item \textbf{Elitism:} The combined parent--offspring pool of size
          $2N$ is reduced back to $N$ by fronts and crowding distance.
\end{enumerate}

\noindent
NSGA-II configuration summary:
\begin{itemize}
    \item population size $N = 80$, maximum generations = 500,
    \item crossover probability $p_c = 0.9$ (uniform),
    \item mutation probability $p_m = 0.08$ (bit-flip),
    \item early stopping: HV improvement $< 0.0005$ over 40
          consecutive generations (minimum 100 generations run).
\end{itemize}

%-----------------------------------------
\subsection{MOPSO Optimization Layer}
\label{sec:mopso_layer}

MOPSO is deployed as a comparative benchmark. Unlike NSGA-II, MOPSO
maintains an external archive of bounded size ($|\mathcal{A}| = 100$)
storing all historically discovered non-dominated solutions. Each
particle $j$ carries a continuous position vector
$\mathbf{p}_j \in [0,1]^N$, binarised at evaluation time via:
\begin{equation}
x_j^{(k)} =
\begin{cases}
1, & \text{if } p_j^{(k)} > 0.5,\\
0, & \text{otherwise.}
\end{cases}
\label{eq:binarise}
\end{equation}
\textit{where:} $p_j^{(k)}$ is the $k$-th component of particle $j$'s
continuous position vector; the threshold $0.5$ maps the continuous space
to binary station-selection decisions.

\medskip\noindent
The velocity update follows the standard constriction form:
\begin{equation}
\mathbf{v}_j \leftarrow
w\,\mathbf{v}_j
+ c_1\,\mathbf{r}_1 \odot
  \bigl(\mathbf{p}_j^{\mathrm{best}} - \mathbf{p}_j\bigr)
+ c_2\,\mathbf{r}_2 \odot
  \bigl(\mathbf{g}^{\mathrm{best}} - \mathbf{p}_j\bigr),
\label{eq:velocity}
\end{equation}
\textit{where:}
\begin{itemize}
  \item $w = 0.4$ — inertia weight controlling velocity momentum,
  \item $c_1 = c_2 = 1.5$ — cognitive and social acceleration
        coefficients,
  \item $\mathbf{r}_1,\mathbf{r}_2 \sim \mathcal{U}(0,1)^N$ — random
        perturbation vectors drawn independently each iteration,
  \item $\mathbf{p}_j^{\mathrm{best}}$ — personal best position of
        particle $j$ over all past iterations,
  \item $\mathbf{g}^{\mathrm{best}}$ — global guide selected from archive
        $\mathcal{A}$ with probability inversely proportional to local
        crowding distance,
  \item $\odot$ — element-wise (Hadamard) product.
\end{itemize}

\noindent
MOPSO configuration: swarm size = 80, $|\mathcal{A}| = 100$,
$w = 0.4$, $c_1 = c_2 = 1.5$; early stopping: $\Delta\mathrm{HV} <
0.0005$ over 40 iterations.

%-----------------------------------------
\subsection{M/M/$c$/K Queueing Evaluation Layer}
\label{sec:queue_layer}

Service quality at each selected station is evaluated using an
M/M/$c$/K finite-capacity queueing model. For station $j$ with
$c_j$~chargers and $W = 8$ physical waiting bays, the system capacity
is $K = c_j + 8$. The offered traffic load is:
\begin{equation}
a = \frac{\lambda_j}{\mu}, \qquad \mu = 2\ \text{sessions/hr.}
\label{eq:load}
\end{equation}
\textit{where:} $\lambda_j$ — arrival rate at station $j$ (sessions/hr);
$\mu = 2$ sessions/hr — per-charger service rate (mean service time
30~min); $a$ — offered load in Erlangs.

\medskip\noindent
The steady-state empty-system probability is:
\begin{equation}
p_0 =
\left[
\sum_{n=0}^{c_j} \frac{a^n}{n!}
+
\sum_{n=c_j+1}^{K} \frac{a^n}{c_j!\; c_j^{\,n-c_j}}
\right]^{-1}.
\label{eq:p0}
\end{equation}
\textit{where:}
\begin{itemize}
  \item $p_0$ — probability the system is completely empty,
  \item first sum: states $0 \le n \le c_j$ (arrivals served immediately),
  \item second sum: states $c_j < n \le K$ (queue forms behind $c_j$
        busy servers),
  \item $K = c_j + 8$ — maximum system capacity.
\end{itemize}

\medskip\noindent
Blocking probability (system full; incoming vehicle turned away):
\begin{equation}
P_{\mathrm{block}} = p_K
= \frac{a^K}{c_j!\; c_j^{\,K-c_j}} \cdot p_0.
\label{eq:pblock}
\end{equation}
\textit{where:} $p_K$ — steady-state probability of exactly $K$ vehicles
in system; $P_{\mathrm{block}} \leq 0.10$ is the hard constraint —
violation adds $+50$~min penalty to $f_4$ (Eq.~\eqref{eq:f4}).

\medskip\noindent
Mean queue waiting time via Little's Law:
\begin{equation}
W_{q,j}
= \frac{L_{q,j}}{\lambda_j\,(1 - P_{\mathrm{block}})} \times 60
\quad\text{(minutes).}
\label{eq:wq}
\end{equation}
\textit{where:}
\begin{itemize}
  \item $L_{q,j}$ — mean number of vehicles waiting (not in service)
        at station $j$,
  \item $\lambda_j(1-P_{\mathrm{block}})$ — effective admitted throughput
        (vehicles/hr entering system),
  \item $\times 60$ — unit conversion from hours to minutes.
\end{itemize}

\medskip\noindent
The day is segmented into three Time-of-Use periods with independent
arrival rates $\lambda_p$ and demand multipliers:

\begin{center}
\renewcommand{\arraystretch}{1.3}
\begin{tabular}{llcc}
\hline
\textbf{Period} & \textbf{Hours} & \textbf{Multiplier}
                & \textbf{Price (INR/session)}\\
\hline
Peak   & 08--10, 17--19 (3~h)  & $\times 2.8$  & \rupee210 \\
Normal & 10--17, 19--22 (13~h) & $\times 1.0$  & \rupee150 \\
Idle   & 22--06 (8~h)          & $\times 0.15$ & \rupee112 \\
\hline
\end{tabular}
\end{center}

\noindent
The reported $\overline{W}_{q,j}$ is the time-weighted average of
$W_{q,j}$ across all three ToU periods. Any station with
$P_{\mathrm{block}} > 10\%$ incurs a 50-minute penalty on objective
$f_4$, forcing the evolutionary optimizer to provision adequate charger
capacity at demand bottlenecks.

%-----------------------------------------
\subsection{Algorithm Performance Evaluation}
\label{sec:performance_eval}

Both algorithms are compared using three standard multi-objective
performance indicators computed over the four-dimensional objective space.
Let $\mathcal{P}$ denote the approximated Pareto front and $\mathbf{r}$
the reference (nadir) point set at 10\% above the empirical nadir of the
combined fronts.

\subsubsection*{Hypervolume (HV)}
\begin{equation}
\mathrm{HV}(\mathcal{P},\mathbf{r})
= \lambda_4\!\left(
    \bigcup_{\mathbf{s}\in\mathcal{P}}
    [\mathbf{s},\mathbf{r}]
  \right).
\label{eq:hv}
\end{equation}
\textit{where:}
\begin{itemize}
  \item $\mathcal{P}$ — approximated Pareto front produced by the
        algorithm,
  \item $\mathbf{r} \in \mathbb{R}^4$ — reference nadir point (set at
        10\% above empirical nadir of combined fronts),
  \item $\lambda_4(\cdot)$ — four-dimensional Lebesgue (hyper-volume)
        measure,
  \item $[\mathbf{s},\mathbf{r}]$ — hyper-rectangle dominated by
        solution $\mathbf{s}$ and bounded by $\mathbf{r}$.
\end{itemize}
Higher HV indicates a larger dominated objective space — preferred.

\subsubsection*{Generational Distance (GD)}
\begin{equation}
\mathrm{GD}(\mathcal{P},\mathcal{P}^*)
= \frac{1}{|\mathcal{P}|}
  \sum_{\mathbf{s}\in\mathcal{P}}
  \min_{\mathbf{s}^*\in\mathcal{P}^*}
  \|\mathbf{s} - \mathbf{s}^*\|_2.
\label{eq:gd}
\end{equation}
\textit{where:}
\begin{itemize}
  \item $\mathcal{P}^*$ — reference Pareto front (union of solutions
        from both algorithms),
  \item $\|\mathbf{s} - \mathbf{s}^*\|_2$ — Euclidean distance in
        $\mathbb{R}^4$ from approximated solution $\mathbf{s}$ to
        nearest reference point $\mathbf{s}^*$.
\end{itemize}
Lower GD indicates closer proximity to the true front — preferred.

\subsubsection*{Spread ($\Delta$)}
\begin{equation}
\Delta =
\frac{d_f + d_l
      + \displaystyle\sum_{i=1}^{|\mathcal{P}|-1}|d_i - \bar{d}\,|}
     {d_f + d_l + (|\mathcal{P}|-1)\,\bar{d}}.
\label{eq:spread}
\end{equation}
\textit{where:}
\begin{itemize}
  \item $d_i$ — crowding distance between the $i$-th and $(i{+}1)$-th
        consecutive Pareto solutions along the front,
  \item $\bar{d}$ — mean crowding distance across all solution pairs,
  \item $d_f$, $d_l$ — distances from the first/last solutions to the
        extreme boundary points of the front.
\end{itemize}
Lower $\Delta$ indicates a more uniformly distributed Pareto set — preferred.

%-----------------------------------------
\section{Summary}
\label{sec:arch_summary}

\noindent
The proposed architecture unifies four-objective Pareto optimisation,
stochastic demand modelling, finite-buffer queueing analysis, and formal
algorithm comparison into a single computational framework.
Table~\ref{tab:arch_summary} summarises the key design choices, their
justifications, and corresponding equation cross-references.

\begin{table}[htbp]
\centering
\small
\caption{Summary of key architectural design choices and their justifications.}
\label{tab:arch_summary}
\setlength{\tabcolsep}{5pt}
\renewcommand{\arraystretch}{1.3}
\begin{tabular}{|p{9em}|p{10.5em}|p{12em}|p{4em}|}
\hline
\textbf{Component} & \textbf{Design Choice} & \textbf{Justification}
                   & \textbf{Eq.} \\
\hline
Demand model &
Logistic S-curve + CAGR + Monte~Carlo &
Captures technology adoption lifecycle and stochastic variability &
\eqref{eq:cagr}--\eqref{eq:mc} \\
\hline
Objective~3 &
Net ROI (profit\,/\,CAPEX) instead of raw profit &
Prevents degeneracy with coverage objective ($r > 0.99$) &
\eqref{eq:f3} \\
\hline
Queue model &
M/M/$c$/K with $K = c + 8$ &
Finite waiting bays more realistic than Erlang-B loss model &
\eqref{eq:p0}--\eqref{eq:wq} \\
\hline
Demand segmentation &
Three ToU periods (Peak\,/\,Normal\,/\,Idle) &
Captures bimodal daily arrival pattern in urban EV charging &
--- \\
\hline
Blocking constraint &
$P_{\text{block}} \leq 10\%$ as hard penalty in loop &
Forces optimizer to provision adequate capacity at bottlenecks &
\eqref{eq:pblock} \\
\hline
Algorithms &
NSGA-II and MOPSO in parallel &
Complementary strengths; enables formal comparative evaluation &
\eqref{eq:binarise}--\eqref{eq:velocity} \\
\hline
Performance metrics &
HV, GD, Spread &
Standard MOEA indicators; algorithm-independent comparison &
\eqref{eq:hv}--\eqref{eq:spread} \\
\hline
\end{tabular}
\end{table}
